\documentclass[aspectratio=169]{beamer}

\usetheme{metropolis}
\setbeamertemplate{navigation symbols}{}

\usepackage[T1]{fontenc}
\usepackage[utf8]{inputenc}
\usepackage{amsmath,amssymb}
\usepackage{booktabs}
\usepackage{graphicx}
\usepackage{array}

\title{Sparse Autoencoders (SAE) for LLM Interpretability}
\subtitle{From superposition to monosemantic features}
\author{330 ML Lab}
\date{}

\newcommand{\slideimg}[1]{%
  \begin{center}
    \includegraphics[width=0.88\textwidth,height=0.58\textheight,keepaspectratio]{#1}
  \end{center}
}

\begin{document}

\begin{frame}
  \titlepage
\end{frame}

\begin{frame}{Journey to SAE-based Model Interpretation}
\begin{enumerate}
  \item Understanding the superposition problem in LLMs
  \item Collecting intuition on how SAE can help
  \item Data collection and training of SAE variants
  \item Architectural shift from vanilla SAE to TopK and FC\_SAE
  \item Training statistics and health checks
  \item Results on reconstruction quality and interpretability
  \item Interpretation of learned features and failure cases
  \item Inference UI walkthrough for practical feature analysis
\end{enumerate}
\end{frame}

\begin{frame}{1) Superposition in LLMs}
\textbf{Core issue}
\begin{itemize}
  \item LLM activations encode more concepts than available neurons
  \item A single neuron often responds to multiple unrelated patterns
  \item This is observed as \textbf{polysemanticity}
\end{itemize}

\vspace{0.4em}
\textbf{Consequence}
\begin{itemize}
  \item Neuron-level interpretation becomes unstable and ambiguous
  \item Hard to assign one clean concept to one neuron
\end{itemize}
\end{frame}

\begin{frame}{Superposition Intuition}
\textbf{Geometric view}
\begin{itemize}
  \item Suppose latent space has dimension $d$
  \item Model needs to represent $m$ concepts with $m \gg d$
  \item Concepts are packed as overlapping linear combinations
\end{itemize}

So observed activation vector is a mixture:
\[
z \approx \sum_{i=1}^{m} a_i f_i
\]

\begin{itemize}
  \item $f_i$: underlying feature directions
  \item $a_i$: feature strengths
\end{itemize}
\end{frame}

\begin{frame}{Why This Hurts Interpretability}
\begin{itemize}
  \item Raw neuron activations are not aligned with human concepts
  \item A neuron can fire for different causes in different contexts
  \item Causal interventions at neuron level can be noisy
\end{itemize}

\vspace{0.5em}
\textbf{Goal of SAE-based MI}
\begin{itemize}
  \item Decompose dense/polysemantic representations
  \item Recover sparse, more monosemantic feature activations
\end{itemize}
\end{frame}

\begin{frame}{2) How Sparsity Solves Superposition}
\textbf{Key idea: sparse dictionary coding}

Learn:
\begin{itemize}
  \item an overcomplete feature dictionary ($m > d$)
  \item sparse coefficients so only a few features are active per token
\end{itemize}

Then each input is reconstructed by a small active subset:
\[
z \rightarrow h(z) \rightarrow \hat{z}
\]

\begin{itemize}
  \item Sparsity encourages feature specialization
  \item Specialization improves interpretability
\end{itemize}
\end{frame}

\begin{frame}{Sparse Representation Principle}
Given encoder output $h(z)$:
\begin{itemize}
  \item Most entries in $h(z)$ should be 0 (or near 0)
  \item Only informative features should activate
\end{itemize}

This pushes the model away from:
\begin{itemize}
  \item ``everything somewhat active''
\end{itemize}

Toward:
\begin{itemize}
  \item ``few meaningful features strongly active''
\end{itemize}

Result: cleaner feature-to-concept mapping.
\end{frame}

\begin{frame}{3) SAE Intro: Architecture}
Standard SAE components:
\begin{enumerate}
  \item \textbf{Encoder} maps LLM activation $z \in \mathbb{R}^d$ to sparse code $h(z) \in \mathbb{R}^m$
  \item \textbf{Decoder} reconstructs $\hat{z}$ from sparse code
\end{enumerate}

Typical form (survey notation):
\[
\hat{z} = h(z) W_{dec} + b_{dec}
\]

\begin{itemize}
  \item Overcomplete setting: $m \gg d$
  \item Decoder columns often interpreted as learned features
\end{itemize}
\end{frame}

\begin{frame}{SAE Loss (L1-Regularized Form)}
From the survey's standard objective:
\[
L(z) = \|z - \hat{z}\|_2^2 + \alpha \|h(z)\|_1
\]

\begin{itemize}
  \item Reconstruction term keeps information fidelity
  \item $L_1$ term enforces sparse activations
  \item $\alpha$ controls trade-off
\end{itemize}

\textbf{Trade-off}
\begin{itemize}
  \item High $\alpha$: more sparse, possibly worse reconstruction
  \item Low $\alpha$: better reconstruction, less disentanglement
\end{itemize}
\end{frame}

\begin{frame}{SAE Loss (KL Sparsity Form)}
From Andrew Ng sparse AE notes:
\[
J_{sparse}(W,b) = J(W,b) + \beta \sum_{j=1}^{s_2} KL(\rho\,\|\,\hat{\rho}_j)
\]

Where:
\begin{itemize}
  \item $\rho$: target average activation (small, e.g., 0.05)
  \item $\hat{\rho}_j$: empirical average activation of hidden unit $j$
  \item $\beta$: sparsity penalty weight
\end{itemize}

Interpretation:
\begin{itemize}
  \item Explicitly forces each hidden unit to be rarely active on average
\end{itemize}
\end{frame}

\begin{frame}{Tuning Parameters (Practical)}
\textbf{Core knobs}
\begin{itemize}
  \item Dictionary size $m$ (expansion factor)
  \item Sparsity strength: $\alpha$ (L1) or $(\rho, \beta)$ (KL)
  \item Learning rate and optimizer settings
  \item Batch size and token sampling strategy
\end{itemize}

\textbf{Typical tuning workflow}
\begin{enumerate}
  \item Start with stable reconstruction
  \item Increase sparsity gradually
  \item Track dead features and feature usage spread
  \item Select checkpoint with best interpretability/reconstruction balance
\end{enumerate}
\end{frame}

\begin{frame}{TopK SAE}
TopK SAE enforces sparsity by selection, not by $L_1$ penalty.

Encoder idea:
\[
\tilde{h}(z) = TopK\big(W_{enc}(z - b_{pre})\big)
\]

Training objective (survey form):
\[
L(z) = \|z - \hat{z}\|_2^2 + \alpha L_{aux}
\]

\begin{itemize}
  \item Keep only top-$K$ activations per token
  \item Avoids direct magnitude shrinkage from L1
  \item Auxiliary loss helps stabilization / dead-latent reduction
\end{itemize}
\end{frame}

\begin{frame}{FC\_SAE (Feature Choice SAE)}
Feature Choice SAE changes sparsity allocation direction.

Instead of:
\begin{itemize}
  \item fixed active features per token
\end{itemize}

It enforces:
\begin{itemize}
  \item each feature gets assigned to a controlled number of tokens
\end{itemize}

Survey statement (Ayonrinde, 2024):
\[
\sum_j S_{i,j} = m,\ \forall i
\]

\begin{itemize}
  \item $S_{i,j}$ indicates whether feature $i$ is active for token $j$
  \item Promotes uniform feature utilization
  \item Targets dead-feature collapse directly
\end{itemize}
\end{frame}

\begin{frame}{4) Architectural Shift: Why Move Beyond Vanilla SAE?}
Reported issues in standard SAE variants:
\begin{itemize}
  \item Shrinkage bias from $L_1$ penalties
  \item Dead latents (features that never activate)
  \item Rigid sparsity allocation
\end{itemize}

New architectures aim to keep:
\begin{itemize}
  \item strong reconstruction
  \item sparse features
  \item healthier feature utilization
\end{itemize}
\end{frame}

\begin{frame}{Architectural Comparison (Quick)}
\small
\begin{tabular}{>{\raggedright\arraybackslash}p{0.2\textwidth} >{\raggedright\arraybackslash}p{0.3\textwidth} >{\raggedright\arraybackslash}p{0.22\textwidth} >{\raggedright\arraybackslash}p{0.2\textwidth}}
\toprule
Variant & Sparsity mechanism & Strength & Main concern \\
\midrule
Vanilla SAE & $L_1$ on activations & Simple and standard & Shrinkage bias \\
TopK SAE & Hard top-$K$ per token & Direct control, cleaner sparsity & Choosing good $K$ \\
FC\_SAE & Token allocation per feature & Better feature usage & More complex assignment dynamics \\
\bottomrule
\end{tabular}
\end{frame}

\begin{frame}{5) Training Statistics -- Big Picture}
\slideimg{stats_training_curve.png}
\begin{itemize}
  \item Training settles into a stable pattern instead of bouncing around wildly
  \item Early epochs do most of the heavy lifting, later epochs are mostly polish
  \item This is the point where we can start tuning for interpretability, not just reconstruction
\end{itemize}
\end{frame}

\begin{frame}{5a) Training Health Check}
\slideimg{stats_dead_feature_ratio.png}
\begin{itemize}
  \item Dead-feature ratio is the ``are we wasting neurons?'' dashboard
  \item The trend suggests feature usage is becoming healthier over time
  \item If this shot up too much, we'd dial back sparsity or adjust learning rate
\end{itemize}
\end{frame}

\begin{frame}{5b) Feature Density Snapshot}
\slideimg{stats_feature_density_hist.png}
\begin{itemize}
  \item Most activations stay near zero, which is exactly what sparse coding wants
  \item A smaller set of features carries strong signal when needed
  \item In plain terms: fewer features talk, but when they do, they say something useful
\end{itemize}
\end{frame}

\begin{frame}{5c) Full Training Curves}
\slideimg{training_history.png}
\begin{itemize}
  \item This gives the complete training-story timeline in one view
  \item Curves are smooth enough to trust the run, not just one lucky batch
  \item Good checkpoint zone: where reconstruction is stable and feature usage looks balanced
\end{itemize}
\end{frame}

\begin{frame}{6) Results -- Reconstruction Quality}
\slideimg{results_reconstruction_compare.png}
\begin{itemize}
  \item Reconstructed activations track the original signal closely
  \item The model keeps key information while still enforcing sparsity
  \item Nice practical takeaway: we are not paying a huge quality tax for interpretability
\end{itemize}
\end{frame}

\begin{frame}{6a) Results -- Variant Comparison}
\slideimg{results_variant_comparison.png}
\begin{itemize}
  \item Vanilla SAE, TopK SAE, and FC\_SAE each land at different trade-off points
  \item TopK/FC-style designs look more consistent on feature usage behavior
  \item If we care about cleaner features, these newer variants are worth the extra complexity
\end{itemize}
\end{frame}

\begin{frame}{6b) Results -- Case-Level Example}
\slideimg{results_case_study_tokens.png}
\begin{itemize}
  \item Token-level view makes differences easier to trust than aggregate metrics alone
  \item We can see where sparsity helps and where reconstruction still misses details
  \item This is the kind of slide that makes results feel concrete in discussion
\end{itemize}
\end{frame}

\begin{frame}{7) Interpretation -- Feature Families}
\slideimg{interpretation_feature_family_overview.png}
\begin{itemize}
  \item Features naturally cluster into themes (syntax, numbers, legal, emphasis, etc.)
  \item This is where monosemantic behavior starts to become visible
  \item Labeling gets much easier once these families emerge consistently
\end{itemize}
\end{frame}

\begin{frame}{7a) Interpretation -- Failure Case}
\slideimg{interpretation_polysemantic_failure.png}
\begin{itemize}
  \item Not every feature is clean; some still mix multiple concepts
  \item These are useful warning signs, not bad news
  \item They tell us where to retune sparsity or inspect data coverage
\end{itemize}
\end{frame}

\begin{frame}{7b) Interpretation -- Activation Evidence}
\slideimg{interpretation_activation_examples.png}
\begin{itemize}
  \item Matched examples help verify whether a label is actually grounded
  \item If examples are coherent, confidence goes up quickly
  \item If examples are messy, we treat the feature as ambiguous and keep investigating
\end{itemize}
\end{frame}

\begin{frame}{8) Inference UI -- Overview}
The inference UI (by ANTLR) has \textbf{three main tabs}:

\vspace{0.4em}
\begin{tabular}{>{\raggedright\arraybackslash}p{0.24\textwidth} >{\raggedright\arraybackslash}p{0.68\textwidth}}
\toprule
Tab & Purpose \\
\midrule
\textbf{SAE Analysis} & Input text -> see top-K activated features per token \\
\textbf{Feature Map} & Search/label features across the full dataset \\
\textbf{Feature Trace} & Trace a specific feature ID over any input text \\
\bottomrule
\end{tabular}

\vspace{0.6em}
Built on top of GPT-2 SAE checkpoints (e.g. \texttt{best\_model\_fc\_sae\_m32\_layer10.pt})
\end{frame}

\begin{frame}{8a) SAE Analysis -- Token-level Feature Inspection}
\slideimg{ui_token_view.jpg}
\begin{itemize}
  \item User pastes any input text and clicks \textbf{Analyze}
  \item Hovering/tapping a token shows its \textbf{top-K activated features} on the right
  \item Here, token \texttt{language} -> top feature \texttt{\#3000: Language Concept} (activation \textbf{13.846})
  \item Remaining features (\#2071, \#3296, ...) show lower but non-trivial activations
\end{itemize}
\end{frame}

\begin{frame}{8b) Feature Map -- Dataset Activation Search}
\slideimg{ui_feature_map.jpg}
\begin{itemize}
  \item Given a \textbf{Feature ID}, scans the dataset (MLCommons/peoples\_speech) for activating sentences
  \item Feature \#6068 labelled \emph{"Emphasis or extent"} -- found 4 matches in 200 sentences
  \item Highlighted token per sentence shows \textbf{exactly which token triggered} the feature
  \item Supports \textbf{Bulk Label} mode to label many features at once via LLM
\end{itemize}
\end{frame}

\begin{frame}{8c) Feature Analysis -- LLM Auto-labeling}
\slideimg{ui_sae_analysis.jpg}
\begin{itemize}
  \item For any feature, the system sends top-activating examples to an LLM for \textbf{automatic labeling}
  \item Feature \#3425 -> labelled \emph{"Legal Court Context"} with confidence \textbf{HIGH}
  \item Top token: \texttt{court} -- consistent across all matched sentences
  \item Parameters: sentences to scan, examples sent to LLM, min activation threshold
\end{itemize}
\end{frame}

\begin{frame}{8d) Feature Analysis -- Activation Examples}
\slideimg{ui_feature_examples.jpg}
\begin{itemize}
  \item Matched sentences displayed with \texttt{>>>token<<<} notation marking the activating token
  \item Each example has a \textbf{prompt-aligned activation score} (e.g. 15.10, 14.58, ...)
  \item Enables manual verification of the LLM-assigned label
  \item Useful for catching mislabeled or polysemantic features
\end{itemize}
\end{frame}

\begin{frame}{8e) Feature Trace -- Single-feature Token Heatmap}
\slideimg{ui_feature_trace.jpg}
\begin{itemize}
  \item Input text + Feature ID -> highlights \textbf{which tokens} activate that feature
  \item Feature \#2873 (\emph{"Numerical Quantity"}) fires on token \texttt{hundred} in the sentence
  \item Right panel: \textbf{Token Activation Map} with underline intensity as a visual cue
  \item Stats shown: active tokens, token count, max activation, threshold
\end{itemize}
\end{frame}

\begin{frame}{Closing}
\begin{itemize}
  \item Superposition motivates feature-level decomposition
  \item Sparsity makes decomposition practical and interpretable
  \item SAE family evolves from L1-based to allocation-aware designs
  \item TopK SAE and FC\_SAE reflect this architectural shift
\end{itemize}
\end{frame}

\begin{frame}{References}
\begin{itemize}
  \item To be added
\end{itemize}
\end{frame}

\begin{frame}[plain]
\centering
\Large Thank you.
\end{frame}

\end{document}
